% !Mode:: "TeX:UTF-8"

\chapter{总结与展望}
\label{chap:conclusion}

\section{全文工作总结}

本文针对现代实时渲染中光线追踪计算开销大、低采样率下噪声泛滥的问题，基于 Vulkan 图形 API 从零构建了一套高性能的实时混合渲染引擎。全文的主要工作总结如下：

\begin{enumerate}
    \item \textbf{设计并实现了数据驱动的 \texttt{RenderGraph} 调度架构}：
    利用有向无环图（DAG）实现了渲染任务的逻辑解耦，通过自动化的资源生命周期分析与屏障（Barrier）推导，解决了 Vulkan 手动同步繁琐且易错的问题。此外，系统实现了鲁棒的时域历史资源持久化机制，为后续 SVGF 与 TAA 算法的高效执行提供了稳定的信号回溯基础。
    
    \item \textbf{构建了PBR 材质模型的混合渲染管线}：
    整合了光栅化 G-Buffer 提取与硬件加速光线追踪（RT Core）的优势。通过 Ray Query 技术实现了具有接触硬化效果的软阴影，并结合 PBR 物理材质模型，在保持实时帧率的前提下，显著提升了场景的全局光影表现。
    
    \item \textbf{落地了高性能的时空信号重建系统}：
    深入研究并实现了 SVGF 时空方差引导滤波算法，在 1 SPP 的低采样率下通过时域累积与 À-Trous 空间滤波，实现了高质量的降噪效果。同时，集成了 TAA 时域抗锯齿方案，通过色彩裁剪与自适应重投影技术，消除了几何边缘走样，优化画面表现。
\end{enumerate}

\section{研究不足与未来展望}

尽管本文构建的混合渲染系统已具备较高的工程完备性，但在算法深度与功能覆盖上仍存在提升空间：

\begin{enumerate}
    \item \textbf{多弹射全局光照（Multi-bounce GI）}：
    目前的系统主要聚焦于直接光阴影与单次弹射的间接光。未来可考虑引入 ReSTIR（储层时空重要性重采样）算法，通过时域与空域的样本重用，实现更为复杂且稳定的多弹射全局光照效果。
    
    \item \textbf{复杂的材质表现}：
    当前引擎主要支持标准的 Cook-Torrance PBR 材质。未来可进一步扩展对各向异性、次表面散射以及薄透镜折射等高级材质特性的硬件光追支持。
    
    \item \textbf{深度学习超分辨率与光追重建技术}：
    虽然 TAA 解决了大部分抗锯齿问题，但在极端动态下仍有模糊感。集成 NVIDIA DLSS 3.5 (Ray Reconstruction) 或 AMD FSR 等技术，利用 AI 辅助特征提取，将是进一步提升画面的重要方向。

    \item \textbf{\texttt{RenderGraph} 显存物理别名化与后处理增强}：
    目前的 \texttt{RenderGraph} 虽实现了逻辑解耦，但尚未完全挖掘显存物理复用的潜力。未来可引入基于资源生存区间的 Aliasing 策略以进一步压低显存峰值。同时，集成 ACES Filmic 等电影级色调映射曲线，将大幅提升系统的视觉表现力。

\end{enumerate}

总之，实时混合渲染技术正处于从“近似模拟”向“物理真实”跨越的关键期。本文所构建的底层架构为后续探索更为前沿的图形学算法提供了坚实的实验平台。
